\addcontentsline{toc}{section}{چکیده انگلیسی}
\thispagestyle{empty}

\begin{latin}
\begin{center}

{\huge Design and Implementation of a Multi-Agent Digital Twin Framework for Intelligent Monitoring, Prediction, and Decision-Making in Industrial Systems}

\vspace{1cm}

{\LARGE{Navid Nasiri}}

\vspace{0.2cm}

{\small navidnasyry@gmail.com}

\vspace{0.5cm}

2026

\vspace{0.5cm}

Department of Electrical and Computer Engineering

\vspace{0.2cm}

Isfahan University of Technology, Isfahan 84156-83111, Iran

\vspace{0.2cm}

Degree: M.Sc. \hspace*{3cm} Language: Farsi

\vspace{1cm}

{\small\textbf{Supervisor: Prof. Masoudreza Hashemi (hashemim@iut.ac.ir)}}

{\small\textbf{Advisor: Prof. Aref Karimi Afshar (a.karimiafshar@iut.ac.ir)}}
\end{center}
~\vfill


\noindent\textbf{Abstract}

\begin{small}
\baselineskip=0.6cm

With the increasing complexity of industrial equipment and the costs associated with unplanned downtime, developing systems for continuous condition monitoring, failure prediction, and appropriate decision execution has become essential. Despite the capability of digital twins in connecting physical assets to their virtual representations, integrating remaining useful life prediction, near-real-time monitoring, human–system interaction, and safe decision-making within scalable architectures remains a significant challenge. This research aims to design, implement, and evaluate a multi-agent digital twin framework for industrial equipment monitoring and failure prediction, leveraging deep learning and large language models. The proposed architecture consists of five layers: physical entity, data acquisition, data integration, digital twin agents, and application agents. To evaluate the framework, a digital twin of turbofan engines was developed based on the FD001 subset of NASA’s CMAPSS dataset, comprising run-to-failure data from 100 training engines and 100 test engines. To investigate the effect of loss function on prediction behavior, two configurations of a three-layer LSTM-based architecture were evaluated. The baseline model, trained with MSE loss, achieved an RMSE of 15.93 and an MAE of 11.95. The second model employed an asymmetric Pinball loss function with a quantile parameter of τ = 0.125, yielding an RMSE of 20.49 and an MAE of 14.72. Despite the increase in numerical error, the under-estimation ratio improved from 0.51 to 0.71, demonstrating that the asymmetric loss function successfully shifted the prediction bias toward under-estimation, albeit at the cost of higher overall error. Independent monitoring, prediction, status dashboard, and language interaction agents were also implemented to support near-real-time twin state updates, delivery of required operational information, and execution of predefined operational commands. The results indicate that the proposed framework provides a modular, interchangeable, and extensible platform for safe, interactive, and adaptable predictive maintenance, suitable for integration with other industrial assets and datasets.


\end{small}

\vspace{0.5 cm}

\noindent \textbf{Key Words}: Digital Twin, Multi-Agent System, Multi-Agent Digital Twin, Large Language Model, Predictive Maintenance, Remaining Useful Life, LSTM, NASA CMAPSS, Real-Time Monitoring, Explainable Artificial Intelligence.

\end{latin}